\documentclass[11pt]{article}

\usepackage[utf8]{inputenc}
\usepackage[T1]{fontenc}
\usepackage[spanish]{babel}
\usepackage{amsmath,amssymb,mathtools,bm,graphicx,xcolor}
\usepackage{svg}
\svgsetup{inkscapelatex=false}
\usepackage[a4paper,margin=2.5cm]{geometry}
\usepackage[
    colorlinks=true,
    linkcolor=red,
    citecolor=green,
    urlcolor=red
]{hyperref}

\spanishdecimal{.}

\setlength{\parindent}{0pt}
\setlength{\parskip}{0.45em}
\setlength{\abovedisplayskip}{6pt}
\setlength{\belowdisplayskip}{6pt}
\setlength{\abovedisplayshortskip}{4pt}
\setlength{\belowdisplayshortskip}{4pt}

% Formato de encabezado y problemas
\newcommand{\encabezado}[3]{%
{\Large\bfseries #1\par}
\vspace{2pt}
{\normalsize #2\hfill #3\par}
\vspace{6pt}\hrule\vspace{8pt}}

\newcommand{\problema}[2]{%
\vspace{0.55em}
\noindent\textbf{#1.}\enspace{\small #2}\par
\vspace{0.3em}}

\newcommand{\inciso}[1]{%
\vspace{0.35em}
\noindent\textbf{(#1)}\enspace}

% Comandos auxiliares
\newcommand{\dd}{\,\mathrm{d}}
\newcommand{\ii}{\mathrm{i}}
\newcommand{\e}{\mathrm{e}}
\newcommand{\sen}{\operatorname{sen}}
\newcommand{\grad}{\nabla}
\newcommand{\laplaciano}{\nabla^2}
\newcommand{\R}{\mathbb{R}}

\begin{document}

\encabezado{Problema de Dirichlet en una esfera}{Física Matemática II}{J.R., G.R.}

\problema{1}{Determine el potencial $\phi(r,\theta,\varphi)$, finito dentro de una esfera de radio $R$, de modo que satisfaga la ecuación de Laplace y la condición de borde
\begin{equation}
\label{cdb-dirichlet-esfera}
\phi(R,\theta,\varphi)=\cos\theta-3\sen^2\theta.
\end{equation}}

\textbf{Solución:} Observamos primero que la condición de borde (\ref{cdb-dirichlet-esfera}) no depende del ángulo azimutal $\varphi$. Por lo tanto el problema tiene simetría axial y podemos buscar una solución de la forma
\begin{equation*}
\phi=\phi(r,\theta).
\end{equation*}
La solución general de la ecuación de Laplace en coordenadas esféricas, finita en el origen, queda dada por
\begin{equation}
\label{solucion-general-interior-dirichlet}
\phi_i(r,\theta)=\sum_{n=0}^{\infty}a_nr^nP_n(\cos\theta),
\qquad r<R.
\end{equation}

Para imponer la condición de borde, definimos
\begin{equation*}
x:=\cos\theta.
\end{equation*}
Entonces, como $\sen^2\theta=1-\cos^2\theta=1-x^2$, se tiene que
\begin{align}
\label{cdb-en-x-dirichlet}
\cos\theta-3\sen^2\theta
&=x-3(1-x^2)\\
&=x-3+3x^2.\nonumber
\end{align}
Ahora queremos escribir esta expresión en términos de polinomios de Legendre. Recordemos que
\begin{equation*}
P_0(x)=1,
\qquad
P_1(x)=x,
\qquad
P_2(x)=\frac{1}{2}(3x^2-1).
\end{equation*}
De la última identidad se obtiene
\begin{equation*}
3x^2-1=2P_2(x)
\quad\Rightarrow\quad
3x^2=2P_2(x)+1.
\end{equation*}
Sustituyendo esto en (\ref{cdb-en-x-dirichlet}), resulta
\begin{align*}
x-3+3x^2
&=P_1(x)-3P_0(x)+2P_2(x)+P_0(x)\\
&=-2P_0(x)+P_1(x)+2P_2(x).
\end{align*}
Por lo tanto, la condición de borde se puede escribir como
\begin{equation}
\label{cdb-legendre-dirichlet}
\phi(R,\theta)=-2P_0(\cos\theta)+P_1(\cos\theta)+2P_2(\cos\theta).
\end{equation}

Evaluamos ahora (\ref{solucion-general-interior-dirichlet}) en $r=R$:
\begin{equation}
\label{interior-evaluado-R-dirichlet}
\phi_i(R,\theta)=\sum_{n=0}^{\infty}a_nR^nP_n(\cos\theta).
\end{equation}
Comparando (\ref{interior-evaluado-R-dirichlet}) con (\ref{cdb-legendre-dirichlet}), y usando que los polinomios $P_n$ son linealmente independientes, encontramos
\begin{equation*}
a_0=-2,
\qquad
 a_1R=1,
\qquad
 a_2R^2=2,
\qquad
 a_n=0 \quad (n\geq 3).
\end{equation*}
Luego,
\begin{equation}
\label{coeficientes-dirichlet}
a_0=-2,
\qquad
 a_1=\frac{1}{R},
\qquad
 a_2=\frac{2}{R^2},
\qquad
 a_n=0 \quad (n\geq 3).
\end{equation}
Sustituyendo (\ref{coeficientes-dirichlet}) en (\ref{solucion-general-interior-dirichlet}), deducimos que
\begin{equation*}
\phi_i(r,\theta)=-2+\frac{r}{R}P_1(\cos\theta)+2\frac{r^2}{R^2}P_2(\cos\theta).
\end{equation*}
Finalmente, usando $P_1(\cos\theta)=\cos\theta$ y $2P_2(\cos\theta)=3\cos^2\theta-1$, obtenemos
\begin{equation*}
\boxed{
\phi_i(r,\theta)=-2+\frac{r}{R}\cos\theta+\frac{r^2}{R^2}\left(3\cos^2\theta-1\right)},
\qquad r<R.
\end{equation*}

Si además se quiere escribir la solución exterior que se anula para $r\to\infty$, consideramos
\begin{equation}
\label{solucion-general-exterior-dirichlet}
\phi_e(r,\theta)=\sum_{n=0}^{\infty}\frac{b_n}{r^{n+1}}P_n(\cos\theta),
\qquad r>R.
\end{equation}
Al imponer la misma condición de borde en $r=R$, tenemos
\begin{equation*}
\frac{b_0}{R}=-2,
\qquad
\frac{b_1}{R^2}=1,
\qquad
\frac{b_2}{R^3}=2,
\end{equation*}
por lo que
\begin{equation*}
b_0=-2R,
\qquad
b_1=R^2,
\qquad
b_2=2R^3.
\end{equation*}
Sustituyendo estos coeficientes en (\ref{solucion-general-exterior-dirichlet}), queda
\begin{equation*}
\boxed{
\phi_e(r,\theta)=-2\frac{R}{r}+\frac{R^2}{r^2}\cos\theta+\frac{R^3}{r^3}\left(3\cos^2\theta-1\right)},
\qquad r>R.
\end{equation*}


\section*{Visualización numérica}

En los mapas angulares se usa la latitud $\lambda=\pi/2-\theta$ para visualizar la esfera como un rectángulo. El color representa el valor del potencial impuesto o calculado. Las líneas negras indican ceros de la función o curvas equipotenciales, según corresponda.

\begin{figure}[h]
\centering
\includesvg[width=0.78\textwidth]{figures/fig_01_condicion_borde}
\caption{Condición de borde $\phi(R,\theta)=\cos\theta-3\sen^2\theta$ y sus contribuciones angulares. Esta descomposición permite identificar qué términos de Legendre aparecen en la solución.}
\end{figure}

\begin{figure}[h]
\centering
\includesvg[width=0.86\textwidth]{figures/fig_05_esfera_coloreada}
\caption{Mapa angular de la condición de borde impuesta sobre la superficie de la esfera. El color representa el valor de $\phi(R,\theta)$ en cada latitud. Las líneas negras indican los ceros de la C.d.B., donde el potencial impuesto cambia de signo.}
\end{figure}

\begin{figure}[h]
\centering
\includesvg[width=0.70\textwidth]{figures/fig_02_potencial_interior}
\caption{Potencial interior en un corte meridiano. El color representa $\phi_i(r,\theta)$ para $r<R$. La circunferencia negra indica la superficie $r=R$ y las líneas delgadas son curvas equipotenciales.}
\end{figure}

\begin{figure}[h]
\centering
\includesvg[width=0.70\textwidth]{figures/fig_03_potencial_exterior}
\caption{Extensión exterior del potencial en un corte meridiano, escogida de modo que se anule cuando $r\to\infty$. El color representa $\phi_e(r,\theta)$ y las curvas delgadas indican equipotenciales.}
\end{figure}

\begin{figure}[h]
\centering
\includesvg[width=0.78\textwidth]{figures/fig_04_descomposicion_frontera}
\caption{Perfiles radiales interiores y exteriores evaluados en direcciones fijas. La línea vertical marca la superficie $r=R$. La figura muestra cómo la solución interior y la exterior coinciden en la frontera, pero tienen distinta dependencia radial.}
\end{figure}

\end{document}
